\documentclass[10pt,a4paper]{amsart}
\usepackage[margin=19mm]{geometry}
\usepackage{amsmath,amssymb,mathtools}
\usepackage[scr=rsfs,cal=euler]{mathalfa}
\usepackage{xfrac}
\usepackage[unicode,hidelinks,bookmarks=false]{hyperref}
\newcommand{\ie}{i.e.\ }
\newcommand{\cf}{cf.\ }
\newcommand{\resp}{respectively\ }
\newcommand{\Subsection}{\subsection}
\providecommand{\nobreakdash}{\nobreak\hbox{-}}
\setlength{\emergencystretch}{2em}
\multlinegap0pt
\usepackage{amssymb}
\usepackage{mathtools}
\usepackage{bm}
\usepackage{xcolor}
\usepackage{csquotes}
\usepackage{dsfont}
\usepackage{algorithm}\usepackage{algpseudocode}
\newtheorem{lemma}{Lemma}
\title{Discretization-free exact recovery in geometric community detection}
\author{Maarten Hoeneveld, Moritz Otto, Rapha\"el Sala}
\date{Source: arXiv:2609.01635v1 (CC BY 4.0)}
\begin{document}
\maketitle

\noindent\emph{Source and licence.} Discretization-free exact recovery in geometric community detection. Version arXiv:2609.01635v1 (CC BY 4.0), distributed by arXiv under the Creative Commons Attribution 4.0 International licence, http://creativecommons.org/licenses/by/4.0/, as recorded in the arXiv licence metadata for this exact version and shown on the abstract page https://arxiv.org/abs/2609.01635v1. Full licence notice: \texttt{licenses/arxiv-2609.01635-CC-BY-4.0.txt}.\par\smallskip

\noindent\textit{Editorial scope.} Counted unit: the article's lemma lem:geometric-split (source0.tex lines 1329--1368). The article's complete original proof of it is delivered with it (source0.tex lines 1370--1459, one \texttt{proof} environment of 90 lines). No mathematics is written, repaired or re-derived: the statement and the proof are the article's own lines, in the article's own notation, and no retained line is substituted or re-flowed.  No further source-local context is needed: every cross-reference of every retained line resolves inside the two delivered spans.  Nothing was omitted from any retained span, so the package carries no omission record.  No retained line contains a citation, so the wrapper carries no bibliography and no bibliography entry.  No retained line contains a figure environment, an image-inclusion command or drawing code, so no image is delivered and the package declares no asset dependency.  Editorial formatting only, in addition: an AMS \texttt{amsart} preamble that restates the article's own declarations and macros used by the retained lines, namely the environments \texttt{lemma} on separate counters, together with the standard packages those lines need.  The retained environments print in this document as Lemma 1 in order of appearance, whereas the article prints the same items with the numbering of its own full document.  The complete native source of the article as distributed in the arXiv e-print for this exact version is bundled unchanged as \texttt{sources/209130/source0.tex}.
\begin{lemma}[Geometric split]
\label{lem:geometric-split}
Fix $s\in(0,r]$ and $\delta>0$. There exists
$\gamma=\gamma(s,\delta)>0$ such that, with probability $1-o(1)$,
the following holds simultaneously for every $a\in Z$ and every
labeling
\[
\sigma:V\cap B_{\mathrm{init}}\to Z.
\]

Suppose that there exist $b,b'\in Z$, with $b\neq b'$, such that
\[
\left|
\{u\in V_a\cap B_{\mathrm{init}}:\sigma(u)=b\}
\right|
\geq\delta\log n
\]
and
\[
\left|
\{u\in V_a\cap B_{\mathrm{init}}:\sigma(u)=b'\}
\right|
\geq\delta\log n.
\]
Then there exists a collection $\mathcal E$ of unordered pairs
$\{u,v\}\subseteq V_a\cap B_{\mathrm{init}}$ such that
\[
|\mathcal E|
\geq
\gamma(\log n)^2,
\]
and, for every $\{u,v\}\in\mathcal E$,
\[
\sigma(u)\neq\sigma(v)
\qquad\text{and}\qquad
\|u-v\|
\leq
s(\log n)^{1/d}.
\]
\end{lemma}

\begin{proof}
Rescale $B_{\mathrm{init}}$ by $(\log n)^{-1/d}$, so that it
becomes the fixed cube $[-r,r]^d$.

Choose an integer $m\geq1$ sufficiently large that, if
$[-r,r]^d$ is partitioned into $m^d$ congruent cubes, then any two
points lying either in the same cube or in two face-adjacent cubes
are at Euclidean distance at most $s$.

Let
\[
\mathcal C=\{C_1,\ldots,C_M\},
\qquad M=m^d,
\]
denote the corresponding partition of $B_{\mathrm{init}}$, scaled
back by $(\log n)^{1/d}$.

Since the points of community $a$ form a Poisson point process of
intensity $\lambda\pi_a$, there exists $\alpha>0$ such that, with
probability $1-o(1)$,
\begin{equation}
\label{eq:true-community-cell-occupancy}
|V_a\cap C|
\geq
\alpha\log n
\end{equation}
simultaneously for every $a\in Z$ and $C\in\mathcal C$.
Indeed, there are only $kM=O(1)$ such random variables, and each is
Poisson with mean of order $\log n$.

Work on the event \eqref{eq:true-community-cell-occupancy}. Set
\[
\tau
:=
\min\left\{
\frac{\delta}{2M},
\frac{\alpha}{2k}
\right\}.
\]
For $C\in\mathcal C$ and $j\in Z$, call $j$ \emph{heavy in $C$} if
\[
\left|
\{u\in V_a\cap C:\sigma(u)=j\}
\right|
\geq
\tau\log n.
\]

By \eqref{eq:true-community-cell-occupancy} and the pigeonhole
principle, every cell has at least one heavy label. Moreover, since
\[
\left|
\{u\in V_a\cap B_{\mathrm{init}}:\sigma(u)=b\}
\right|
\geq\delta\log n,
\]
there exists a cell in which $b$ is heavy. The same holds for $b'$.

Suppose first that some cell $C$ contains two distinct heavy labels
$j\neq j'$. Then there are at least $tau^2(\log n)^2$ pairs $\{u,v\}\subseteq V_a\cap C$ with $\sigma(u)=j$and $sigma(v)=j'$.
By the choice of the partition,
\[
\|u-v\|
\leq
s(\log n)^{1/d},
\]
and the result follows.

It remains to consider the case in which every cell has a unique
heavy label. Assign to each cell its unique heavy label. The
face-adjacency graph of the cells is connected. Since some cell has
heavy label $b$ and some other cell has heavy label $b'$, there must
exist two face-adjacent cells $C,C'$ whose heavy labels, say $j$ and
$j'$, are distinct.

Consequently, there are at least $\tau^2(\log n)^2$ pairs
$u\in V_a\cap C$, $v\in V_a\cap C'$ such that
\[
\sigma(u)=j,
\qquad
\sigma(v)=j'.
\]
Again, by construction,
\[
\|u-v\|
\leq
s(\log n)^{1/d}.
\]
Thus, the assertion holds with $\gamma:=\tau^2$.
\end{proof}

\end{document}
